\textbf{1. Kiến thức cần nhớ}
\begin{itemize}
\item \textbf{Định nghĩa:} hai mặt phẳng $(P)$ và $(Q)$ song song nếu chúng không có điểm chung, kí hiệu $(P)\parallel(Q)$.
\item \textbf{Điều kiện song song:} nếu $(P)$ chứa hai đường thẳng \emph{cắt nhau} $a$, $b$ và $a\parallel(Q)$, $b\parallel(Q)$ thì $(P)\parallel(Q)$.
\item \textbf{Tính chất:}
\begin{itemize}
\item Qua một điểm nằm ngoài một mặt phẳng có duy nhất một mặt phẳng song song với mặt phẳng đó.
\item Nếu $(P)\parallel(Q)$ thì mọi đường thẳng nằm trong $(P)$ đều song song với $(Q)$: \quad $\boxed{(P)\parallel(Q),\ a\subset(P)\Rightarrow a\parallel(Q)}$.
\item Nếu $(P)\parallel(Q)$ và mặt phẳng $(R)$ cắt $(P)$ thì $(R)$ cũng cắt $(Q)$, và hai giao tuyến song song với nhau: \quad $\boxed{(R)\cap(P)=a,\ (R)\cap(Q)=b\Rightarrow a\parallel b}$.
\item Hai mặt phẳng phân biệt cùng song song với một mặt phẳng thứ ba thì song song với nhau.
\end{itemize}
\item \textbf{Định lí Thalès trong không gian:} ba mặt phẳng đôi một song song chắn trên hai cát tuyến bất kì những đoạn thẳng tương ứng tỉ lệ:
\[ \dfrac{AB}{A'B'}=\dfrac{BC}{B'C'}=\dfrac{CA}{C'A'}. \]
\item \textbf{Hình lăng trụ:} hai đáy là hai đa giác bằng nhau, nằm trên hai mặt phẳng song song; các cạnh bên song song và bằng nhau; các mặt bên là hình bình hành.\\
\textbf{Hình hộp:} là hình lăng trụ có đáy là hình bình hành; sáu mặt là hình bình hành, các cặp mặt đối diện song song.
\end{itemize}

\textbf{2. Phương pháp giải các dạng}

%% DANG 1H4?4-1
\textbf{Dạng 1. Câu hỏi lý thuyết}
\begin{enumerate}
\item Đối chiếu từng mệnh đề với định nghĩa, điều kiện và tính chất ở trên.
\item Với mệnh đề sai, chỉ ra một phản ví dụ trên hình hộp hoặc hình lập phương.
\end{enumerate}
Lưu ý: hai đường thẳng trong $(P)$ cùng song song với $(Q)$ phải \emph{cắt nhau}; nếu chúng song song thì chưa kết luận được $(P)\parallel(Q)$. Hai đường thẳng nằm trong hai mặt phẳng song song có thể chéo nhau.

%% DANG 1H4?4-2
\textbf{Dạng 2. Chứng minh hai mặt phẳng song song}
\begin{enumerate}
\item Trong $(P)$, chọn hai đường thẳng cắt nhau $a$, $b$.
\item Chứng minh $a\parallel(Q)$ và $b\parallel(Q)$, thường bằng cách chỉ ra $a$, $b$ song song với hai đường thẳng nằm trong $(Q)$ (dùng đường trung bình, hình bình hành, định lí Thalès).
\item Kết luận $(P)\parallel(Q)$. Có thể thay bằng cách chứng minh $(P)$ và $(Q)$ cùng song song với một mặt phẳng thứ ba.
\end{enumerate}

%% DANG 1H4?4-3
\textbf{Dạng 3. Chứng minh đường thẳng song song với mặt phẳng}
\begin{enumerate}
\item Tìm một mặt phẳng $(R)$ chứa $a$ sao cho $(R)\parallel(P)$.
\item Áp dụng tính chất: $a\subset(R)$, $(R)\parallel(P)\Rightarrow a\parallel(P)$.
\end{enumerate}
Lưu ý: cách này dùng khi khó tìm được đường thẳng nằm trong $(P)$ song song với $a$.

%% DANG 1H4?4-4
\textbf{Dạng 4. Xác định mặt phẳng đi qua một điểm và song song với một mặt phẳng}
\begin{enumerate}
\item Trong $(P)$ chọn hai đường thẳng cắt nhau $a$, $b$.
\item Qua điểm $M$ dựng $a'\parallel a$, $b'\parallel b$; khi đó $(Q)=(a',b')$ là mặt phẳng cần tìm.
\item Tìm thiết diện: dùng tính chất giao tuyến song song, lần lượt xác định giao tuyến của $(Q)$ với từng mặt của hình chóp hoặc hình lăng trụ.
\end{enumerate}

%% DANG 1H4?4-5
\textbf{Dạng 5. Xác định mặt phẳng chứa đường thẳng (hoặc đi qua hai điểm) và song song với mặt phẳng hoặc đường thẳng cho trước}
\begin{enumerate}
\item Nếu $(Q)$ chứa $a$ và song song với $(P)$: qua một điểm của $a$ dựng thêm đường thẳng $b$ song song với $(P)$, $b$ cắt $a$; khi đó $(Q)=(a,b)$.
\item Nếu $(Q)$ chứa $a$ (hoặc đi qua hai điểm) và song song với đường thẳng $d$: qua một điểm của $a$ dựng $d'\parallel d$, khi đó $(Q)=(a,d')$.
\item Tìm giao tuyến của $(Q)$ với các mặt của hình, dùng giao tuyến song song, rồi nối các đoạn giao tuyến để được thiết diện.
\end{enumerate}
Lưu ý: thiết diện là đa giác khép kín, mỗi cạnh nằm trên một mặt của hình; mỗi mặt chứa nhiều nhất một cạnh của thiết diện.

%% DANG 1H4?4-6
\textbf{Dạng 6. Bài toán tổng hợp}
\begin{enumerate}
\item Phân tích yêu cầu: chứng minh song song, tìm giao tuyến, dựng thiết diện hay tính tỉ số.
\item Kết hợp tính chất giao tuyến song song với đường trung bình, định lí Thalès trong tam giác và trong không gian.
\item Tính tỉ số độ dài, chu vi hoặc nhận dạng thiết diện (hình bình hành, hình thang) dựa trên các cặp cạnh song song vừa tìm được.
\end{enumerate}

%% DANG 1H4?4-7
\textbf{Dạng 7. Toán thực tế áp dụng hai mặt phẳng song song}
\begin{enumerate}
\item Mô hình hóa: sàn nhà, trần nhà, các tầng nhà, mặt kệ, hai đáy của hộp được coi là các mặt phẳng song song.
\item Áp dụng tính chất giao tuyến song song (ví dụ: mép cắt của một tấm kính với hai mặt song song) hoặc định lí Thalès trong không gian (ví dụ: cầu thang, dây néo xuyên qua các tầng) để lập tỉ lệ rồi tính độ dài cần tìm.
\end{enumerate}
Lưu ý: định lí Thalès trong không gian chỉ áp dụng khi ba mặt phẳng \emph{đôi một song song}; các tỉ số phải lấy tương ứng trên cùng một cát tuyến.
